\section{Time Plan}
\begin{itemize}
    \item \textbf{September:} Literature study on quantum measurement, quantum dots, double quantum dots, microwave resonators, and input-output theory. Familiarization with the experimental setup and the hybrid microwave system. Reproduce basic calculations and results from previous work to establish a starting point for the project.

    \item \textbf{October:} Continue the literature study, with particular focus on input-output theory and the coupling between quantum systems and microwave resonators as well as Fisher information. Develop the theoretical description of the microwave detector and investigate how the reflection coefficient depends on the properties of the coupled quantum-dot system. Extend calculations from previous work and begin outlining the project report.

    \item \textbf{November:} Develop and investigate the model of the coupled quantum-dot and microwave system. Study the single quantum-dot system and then introduce the double quantum dot as an effective two-level system. Investigate the Jaynes--Cummings model and the more general quantum Rabi model, including resonant and dispersive coupling. Perform initial numerical simulations and continue writing the report.

    \item \textbf{December:} Focus on the measurement process and the connection between the quantum state and the experimentally measured reflection coefficient. Investigate measurement backaction, relaxation, and decoherence. Analyse the most relevant results, create figures, and prepare the results and theoretical background for the half-time meeting.

    \item \textbf{January:} Half-time meeting. Assess the progress and results obtained during the first semester and determine the focus of the remaining work. Refine the theoretical model and research questions based on the results obtained so far. Begin a more detailed investigation of quantum non-demolition (QND) measurement and the conditions required for it.

    \item \textbf{February:} Perform the main calculations and simulations concerning quantum measurement of the DQD. Investigate the effects of relaxation and decoherence on the measurement and compare resonant and dispersive coupling regimes. Study detector sensitivity and investigate Fisher information as a measure of the information obtained from the microwave signal.

    \item \textbf{March:} Analyse the validity of the approximations used in the theoretical description and determine the physical regimes in which they are applicable. Compare different theoretical descriptions where appropriate, and compare the theoretical results with available experimental data. Finalize the main analysis and create figures presenting the central results.

    \item \textbf{April:} Consolidate and interpret the results. Complete the discussion of measurement backaction, decoherence, QND measurement, detector sensitivity, and the validity of the approximations. Write the remaining sections of the report and revise the report as a whole.

    \item \textbf{May:} Finalize the report, including proofreading, figures, references, and conclusions. Prepare slides for the thesis defence and perform a dry run of the presentation. Thesis defence.
\end{itemize}